\section{Thesis Introduction}
\label{sec:thesis-introduction}

\subsection{From Quantum Teleportation to Observable Teleportation}
The transfer of quantum information is a cornerstone of quantum computing and communication. In their seminal 1993 paper, Bennett et al. introduced quantum teleportation, a remarkable protocol that allows for the perfect transfer of an unknown quantum state from one location to another using a pre-shared entangled pair and a classical communication channel. This process is fundamental to the development of quantum networks and distributed quantum computing, as it enables the transmission of fragile quantum states without physically sending the carrier particle, thereby circumventing decoherence associated with direct transmission.

However, the teleportation of a full quantum state is not always necessary or desirable. In many physical contexts, particularly in condensed matter and many-body systems, one is often more interested in the local properties of a system, such as energy, charge, or spin, which are described by the expectation values of local observables. This observation motivated the development of Quantum Energy Teleportation (QET), a conceptually distinct protocol first proposed by M. Hotta \cite{Hotta2011}. Unlike state teleportation, QET does not transfer a quantum state. Instead, it leverages the entanglement present in the ground state of a many-body system to allow one party, Alice, to perform a local measurement and, via classical communication, enable a distant party, Bob, to extract energy from his local subsystem. Crucially, this is achieved without violating the law of local energy conservation; the energy appearing at Bob's side is balanced by negative energy injected into the quantum vacuum by Alice's measurement. The protocol's viability was recently confirmed through an experimental demonstration on superconducting quantum hardware by K. Ikeda \cite{Ikeda2023}.

The principles of QET are not limited to energy. The protocol can be generalized to teleport the expectation value of any globally conserved quantity that admits a local representation, a framework known as Quantum Observable Teleportation \cite{IkedaTeleportingCharge}. This extension significantly broadens the scope of applications to include the nonlocal manipulation of local charge, current, and other physical observables, providing a powerful new tool for probing and controlling quantum many-body systems.

\subsection{Quantum Key Distribution}
Parallel to the development of quantum information processing has been the growth of quantum cryptography, which aims to provide communication security guaranteed by the laws of physics. The most developed of these technologies is Quantum Key Distribution (QKD), which allows two parties, Alice and Bob, to establish a shared, secret random key that can be used for classical encryption. The security of QKD protocols, such as the pioneering BB84 protocol \cite{Bennett1984}, stems from the fact that any attempt by an eavesdropper, Eve, to gain information about the key necessarily disturbs the quantum states being transmitted, alerting Alice and Bob to her presence. This stands in stark contrast to classical cryptographic schemes, whose security relies on computational assumptions that may be broken by future advances, including the advent of a quantum computer.

\subsection{A New Paradigm: QKD via Charge Teleportation}
The convergence of quantum observable teleportation and quantum cryptography presents an exciting new frontier. It was recently proposed that the QET framework could be adapted to create a novel QKD protocol \cite{QKDbyQET}. The core idea is to encode information not in the state of a transmitted particle, but in the sign of the expectation value of an observable measured by Bob. In this scheme, Alice performs a local measurement on her part of a shared entangled state and communicates her measurement basis and outcome to Bob. However, she secretly makes a classical binary choice: to send her true measurement outcome or its opposite. Bob uses the information he receives to perform a conditional local operation, after which he measures a local observable. The sign of his resulting expectation value depends directly on Alice's secret classical choice, thereby securely transmitting one bit of the key. An eavesdropper cannot determine Alice's choice from the quantum state or the public classical communication alone.

This thesis investigates the implementation of such a QKD protocol using Quantum Charge Teleportation (QCT). We analyze its performance in the context of the Transverse Field Ising Model (TFIM) with two different coupling geometries: a star-interaction model ($H^{(1)}$), where a central Alice qubit is coupled to all others, and a 1D nearest-neighbor model ($H^{(2)}$). Through extensive numerical simulations and realistic quantum circuit simulations using Qiskit, we perform a comparative analysis of charge versus energy teleportation. We evaluate the protocols based on signal strength, symmetry, scalability with system size ($N$), and robustness against various quantum noise channels, including communication errors and decoherence. Our results aim to quantify the practical advantages of using a charge-based observable for QKD and to identify the optimal parameter regimes for its implementation on near-term quantum devices.
